\subsection{Primary-source audit and main-question search}
\label{subsec:source-audit}
The following records the actual reading in this session. ``Not obtained''
means that bibliographic records were opened but the mathematical body was
not available through the routes tried. No theorem number or proof is
attributed to an unopened body.

\begingroup
\small
\setlength{\tabcolsep}{4pt}
\begin{longtable}{@{}p{0.23\textwidth}p{0.35\textwidth}p{0.37\textwidth}@{}}
\toprule
Source & Opened passages or unsuccessful attempts & Bearing on this record \\
\midrule
\endfirsthead
\toprule
Source & Opened passages or unsuccessful attempts & Bearing on this record \\
\midrule
\endhead
\bottomrule
\endfoot
\citet{IsbellMitchell1973}
& Full supplied article, all three pages: Theorem~1, Lemma~1, Theorem~2,
Corollaries~1--2, Theorem~3.
& Exactness criterion and finite-family reduction; weak-terminal exclusion;
announcement of the increasing-injection example. \\

\citet{IsbellMitchell1976}
& Full supplied article: Proposition~1.1; Lemmas~2.1--2.2;
Proposition~2.3; Theorem~3.1; Lemma~3.2; Proposition~3.3;
coefficient remark on p.~296; Theorem~4.1; Corollaries~4.2--4.4;
Proposition~4.5; Corollary~4.6; note added in proof.
& Purity and affine filteredness; I+III; transport of IV to connected
set diagrams; endomorphism, weak-terminal and bounded-hom exclusions.
The note rejects the claimed implication I+V$\Rightarrow$III in the
1975 paper. \\

\citet{Isbell1968}
& Full Cambridge PDF, four pages. Both unnumbered theorems:
\S1, p.~569, and \S2, p.~571.
& A finite-family exactness criterion and the one-object case.
The minimal-period argument was checked directly. The original has no
numbered ``Theorem~1'' or ``Theorem~2''. \\

\citet{Isbell1974}
& Full text not obtained. Official Annals article page and JSTOR record
opened; JSTOR PDF route returned HTML. The Annals page supplied only
general PDF-access information.
& Bibliography and the 1976 article's attribution only. No original
proof or numbered theorem claimed as read. \\

\citet{Oberst1968}
& Full text not obtained. Springer article page opened; its PDF link
returned the subscription article page. Exact-title, GDZ and EuDML
searches did not locate an accessible full text.
& No unverified theorem number used. Husainov's separately opened
Theorem~2.7 is distinguished from Oberst's original article. \\

\citet{GabrielUlmer1971}
& Full PDF from Douglas Ravenel's academic archive. \S15.1;
Example~15.2(c), p.~187, and its proof in \S15.9, pp.~190--191;
also \S\S14.8--14.14 and \S8.9.
& The original finite-cocone connectedness criterion for siftedness.
The proof in \S15.9 is needed because Example~15.2(c) announces it
before proving it. No total-to-exact theorem was identified in these
passages. \\

\citet{AdamekReiterman1975}
& Full text not obtained. The AMS PDF, article page and DOI endpoint
returned HTTP~403. Two ResearchGate article records were opened and
had no full text.
& Bibliography only; no proof content inferred. The criticism in
Isbell--Mitchell 1976 is attributed to that opened source. \\

\citet{AdamekKoubekReiterman1981}
& Supplied complete scan, source images and OCR. Construction B.1,
B.4--B.5, C.1, and the full C.2 proof I--VI, including V.2.1--V.2.3,
were personally checked by the topology agent under the repaired
conditions (9)--(11); details and index corrections are in Part~C.
& Quasi-filtered construction, incidence inequalities and the
non-affine-equalizable triple. DML metadata's first-name error is
avoided: the third author is Jan Reiterman. \\

\citet{Husainov2002}
& Full author-posted text opened on ResearchGate; \S\S2.1, 2.3,
Theorems~2.5--2.7, Problem~2.5 and corresponding references inspected.
& Affine filteredness, linearization of injective-action diagrams,
and exactness in AB5 abelian categories. Problem~2.5 is not treated
as evidence of present unresolved status. \\

\citet{SedaghatjooKhaksari2017}
& Complete arXiv preprint and published PDF; Propositions~2.1, 2.3;
Theorem~3.2; Lemma~3.7; Corollary~3.8; Theorem~3.9;
Proposition~3.12.
& For the one-element right act, all-product flatness, super-flatness
and existence of a left zero are equivalent (Proposition~3.12).
Finite-product flatness is separately treated in Theorem~3.2.
These are ordinary set tensor products. \\

\citet{LaanEtAl2021}
& Complete primary PDF; Definitions~1.1--1.3, Remark~1.6,
Theorem~3.5, Theorem~4.2, Proposition~4.4 and Proposition~5.9.
& Distinguishes pullback, finite-limit and all-limit flatness;
conditions (P), (E) and local cyclicity. No inference from permutation
module acyclicity to flatness of $\Delta\Z$ is supplied. \\

\citet{BulmanFlemingLaan2001}
& Full text not obtained. Author publication list and its public
publisher share link opened; the share link yielded no readable text.
& A foundational act-flatness reference cited by the two opened
papers above; no original theorem quoted as read. \\
\end{longtable}
\endgroup

The proof checks needed for our uses are as follows. In the proof of
Isbell--Mitchell's Theorem~3.1, the coefficients produced by purity can
initially have different targets. Condition I transports their finite
support to a common target; the objectwise equations remain zero.
Proposition~3.3 then uses the union of the supports of two affine arrows
and condition III, rather than applying IV to affine arrows. In
Theorem~4.1 the periods belong to an endomorphism acting on a set; no
invertibility is assumed. The weakly terminal object is used later in
Corollary~4.4. In Gabriel--Ulmer's terminology, ``kofinal'' is final
for our covariant diagrams and ``kofiltrierend'' is filtered. Their
two-object cocone criterion is therefore finality of the diagonal,
which is siftedness and does not assert totality.

For the main implication we searched variants of ``exact colimits'' with
``aspherical'', ``totally aspherical'', ``homotopical'', ``Isbell'',
``permutation modules'' and ``homotopy flat''; also ``homologically flat
acts'', ``linearization of acts'', ``Tor permutation modules flat monoid'',
``homological dimension atomic topos'', and ``smooth coinvariants
oligomorphic''. We used both search engines and opened primary PDFs where
found. Searches under ``totally aspherical'' also returned the different
test-category terminology about aspherical diagonals, which we did not
substitute for the definition in this problem.

The closest exact question found is \citet{Henry2026}, posted on
22 April 2026. It asks which shapes have all set colimits preserved by
$\Set\hookrightarrow\mathrm{Spaces}$; the comments identify the condition
with weakly contractible components of all categories of elements. No
answer was visible on the opened page. This is a retrieval record, not
evidence that the present main question has a particular mathematical
status.

Further opened primary sources were \citet{StanculescuClasses},
Definition~1.1, Propositions~1.2 and~1.7, and Definition~2.8, on
hocolimit-commutation classes; \citet[Theorem~3.5]{Caramello2013}, on
the atomic-topos/Fra\"iss\'e comparison; and
\citet{HarmanSnowden2024}, \S\S1.3.3--1.3.4, 10--11, on
measure-dependent representation categories. The latter use Schwartz
spaces and a completed algebra depending on a measure; their
semisimplicity cannot be transferred to ordinary smooth $G$-modules
without an additional comparison theorem. The Ramsey sources used in
our sufficient condition are
\citet[Definition~2.1, Example~2.1 and \S4]{Masulovic2026},
\citet[Theorem~3.21 and Corollaries~4.2--4.3]{HubickaNesetril2019}, and
\citet[Theorem~15 and its proof]{Hubicka2025}. No novelty is claimed.

Retrieval programs and their logs are included under \texttt{code/}.
The audit records unsuccessful requests as well as successful ones.
